% Example "vanderpol": the problem, the flow and its limit cycle, the
% relaxation oscillations for growing mu, and the convergence, period and
% stiffness tests. Every number comes from data/, written by
% examples/vanderpol/main.f90: run `fpm run --example vanderpol` from the
% project root, then build this file with lualatex in examples/vanderpol/.
\documentclass[11pt,a4paper]{article}
\usepackage[margin=2.2cm]{geometry}
\usepackage{amsmath,booktabs,array}
\usepackage{pgfplots}
\usepgfplotslibrary{groupplots,colormaps}
\usetikzlibrary{arrows.meta}
\pgfplotsset{compat=1.18}
\pagestyle{empty}

\colorlet{m-euler}{black}
\colorlet{m-order2}{orange!85!black}
\colorlet{m-order3}{green!50!black}
\colorlet{m-rk4}{blue!80!black}
\colorlet{m-dp54}{red!75!black}

\begin{document}
\section*{Example: the Van der Pol oscillator}

\subsection*{1. Problem}
From the SciML problem library (\texttt{prob\_ode\_vanderpol}): an
oscillator whose damping is negative for small amplitudes and positive for
large ones, so that every solution but the origin settles on one periodic
orbit, a limit cycle,
\begin{equation*}
  x' = y, \qquad y' = \mu\bigl((1 - x^2)\,y - x\bigr), \qquad (x, y)(0) = (\sqrt 3, 0).
\end{equation*}
In the time $\tau = \sqrt\mu\, t$ this is the classic
$\ddot x - \varepsilon (1 - x^2)\dot x + x = 0$ with $\varepsilon = \sqrt\mu$.
There is no closed form; the reference is a \texttt{dp54} run at a tiny
step. For large $\mu$ the oscillation becomes a relaxation oscillation: slow
drifts along the branches $|x| \in [1, 2]$ and fast jumps between them. On
the slow branch near $x = \pm 2$ the Jacobian
$\bigl(\begin{smallmatrix} 0 & 1 \\ -\mu(2xy + 1) & \mu(1 - x^2)\end{smallmatrix}\bigr)$
has an eigenvalue close to $\mu(1 - x^2) = -3\mu$: the problem becomes stiff.
\begin{center}
\small
% Written by examples/vanderpol/main.f90: do not edit.
\begin{tabular}{l>{\raggedright\arraybackslash}p{12.4cm}}
\toprule
parameter & $\mu = 1$ (equivalently $\varepsilon = \sqrt\mu = 1$) \\
initial condition & $(x, y) = (\sqrt 3, 0)$ \\
reference & \texttt{dp54}, $\Delta t = 10^{-4}$; period $T = 6.6632868593$ \\
convergence & error $\max_i|u_i - u_i^{\mathrm{ref}}|$ at $t = 10.0$, $\Delta t = 0.10, \dots, 0.10/256$ \\
period & mean spacing of upward zero crossings 10 to 15 of $x$; each crossing solved to the method's accuracy \\
observed order & two finest levels above $10^{-11}$; kept if within 0.3 of the next slope \\
stiffness & largest $\Delta t$ keeping $|x| < 10$ on $0 \le t \le 10.0$, $\mu = 1, 3, 10, 30, 100, 300, 1000$ \\
figures & solutions at $\Delta t = 0.25$, $0 \le t \le 20.0$; limit cycles for $\mu = 1, 4, 16, 64$ \\
methods & \texttt{rk3ws}, \texttt{rk4}, \texttt{dp54} \\
interface & plain arrays, \texttt{rk\_array\_scheme\_t} \\
\bottomrule
\end{tabular}

\end{center}

\newpage
\subsection*{2. Solutions}
\begin{center}
\begin{tikzpicture}
  % ---- (a) the flow and its limit cycle ----
  \begin{axis}[
      name=flow,                                           % anchor for the panels on the right
      width=8.2cm, height=12.4cm,                          % large phase portrait
      xmin=-3.2, xmax=3.2, ymin=-3.9, ymax=3.9,            % the window of the direction field
      xlabel={$x$}, ylabel={$y$},                          % phase plane
      title={(a) phase plane, $\mu = 1$},                  % panel title
      colormap/viridis,                                    % perceptually uniform colours
      colorbar, colorbar style={ylabel={$\ln(|f| + 0.1)$}, width=0.25cm}, % speed scale
      axis background/.style={fill=black!2},               % very light background
    ]
    \addplot[
        point meta={ln(\thisrow{speed} + 0.1)},            % colour by speed, log scale
        quiver={u=\thisrow{u}, v=\thisrow{v}, scale arrows=0.24, colored=mapped color},
        -{Stealth[length=2.2pt]}, line width=0.6pt,        % small arrow heads
      ] table {data/field.dat};                            % the direction field, unit arrows
    \foreach \k in {1,...,5} {
      \addplot[white, line width=1.6pt, forget plot] table[x=x, y=y] {data/start_\k.dat}; % halo
      \addplot[black!55, line width=0.7pt, forget plot] table[x=x, y=y] {data/start_\k.dat}; % transients
    }
    \addplot[only marks, mark=*, mark size=1.6pt, black!55]
      coordinates {(0.1,0) (3,3) (-3,-2.5) (-0.5,1) (2.6,-3.2)};                          % their starts
    \addplot[red!80!black, line width=2pt, restrict expr to domain={\thisrow{t}}{12:20}]
      table[x=x, y=y] {data/reference.dat};                % the limit cycle (reference, t >= 12)
  \end{axis}
  % ---- (b) x(t) of the methods ----
  \begin{axis}[
      name=series,                                         % anchor for (c)
      at={($(flow.north east)+(3.3cm,0.2cm)$)}, anchor=north west, % right of (a), past the colour bar
      width=5.9cm, height=6cm,                             % panel
      xmin=0, xmax=20, ymin=-3, ymax=3,                    % time window, amplitude
      ylabel={$x$}, xlabel={$t$},                          % solution against time
      title={(b) $x(t)$, $\Delta t = 0.25$},               % panel title
      axis background/.style={fill=black!2},               % very light background
      legend style={font=\tiny, at={(0.5,-0.32)}, anchor=north, legend columns=4, draw=none}, % below
    ]
    \addplot[black!25, line width=3pt] table[x=t, y=x] {data/reference.dat};              % reference, wide
    \addlegendentry{reference}
    \addplot[m-rk4, thick, dashed] table[x=t, y=x] {data/solution_rk4.dat};
    \addlegendentry{\texttt{rk4}}
  \end{axis}
  % ---- (c) relaxation oscillations for growing mu ----
  \begin{axis}[
      at={($(series.south west)+(0,-2.8cm)$)}, anchor=north west, % below (b) and its legend
      width=5.9cm, height=6cm,                             % panel
      xmin=12, xmax=20, ymin=-2.6, ymax=3.1,               % after the transients
      ylabel={$x$}, xlabel={$t$},                          % waveform
      title={(c) limit cycles, $\mu = 1, 4, 16, 64$},      % panel title
      axis background/.style={fill=black!2},               % very light background
      cycle list={[colors of colormap={0,300,600,900} of viridis]}, % four viridis colours
      legend style={font=\tiny, at={(0.5,-0.32)}, anchor=north, legend columns=4, draw=none}, % below
    ]
    \foreach \k/\m in {1/1, 2/4, 3/16, 4/64} {
      \addplot+[no marks, thick] table[x=t, y=x] {data/cycle_\k.dat};                     % x(t) for this mu
      \addlegendentryexpanded{$\mu = \m$}
    }
  \end{axis}
\end{tikzpicture}
\end{center}
(a) The direction field (arrows of equal length, coloured by the speed
$|f|$), five trajectories (grey) from inside and outside, and the limit cycle
(red) they all approach within a few turns. (b) At $\Delta t = 0.25$
\texttt{rk4} follows the reference at this scale. (c) As $\mu$ grows the cycle becomes a relaxation
oscillation: slow drifts and fast jumps, with a shorter period in $t$.

\newpage
\subsection*{3. Tests: convergence, period, stiffness}

\paragraph{Prediction.} (i) The solution converges at each method's general
order. (ii) So does the period, when each crossing is located to the
method's accuracy. (iii) On the slow branch near $x = \pm 2$ a step is stable
while $h\cdot 3\mu \le h^*$, $h^*$ the method's limit on the negative real
axis (the \texttt{stability} example): the largest stable step should
approach $h^*/(3\mu)$, so $\mu\,\Delta t_{\max} \to h^*/3$. Freezing the
Jacobian all along the cycle (``along the cycle''), at the points where the
flow contracts, gives a second linear estimate.

\paragraph{Result.}
\begin{center}
\small
\setlength{\tabcolsep}{4pt}
% Written by examples/vanderpol/main.f90: do not edit.
\begin{tabular}{lccccccl}
\toprule
& & \multicolumn{2}{c}{observed order} & \multicolumn{4}{c}{stiffness, $\mu\,\Delta t_{\max}$ at $\mu = 1000$} \\
\cmidrule(lr){3-4}\cmidrule(lr){5-8}
name & order & solution & period & measured & $h^*/3$ & along the cycle & first blows up \\
\midrule
\texttt{rk3ws} & 2 & 2.01 & 3.00 & 0.8070 & 0.8376 & 0.8322 & slow branch \\
\texttt{rk4} & 4 & 4.00 & 4.00 & 0.8225 & 0.9284 & 0.9225 & fast jump \\
\texttt{dp54} & 5 & --- & 5.98 & 0.8469 & 1.1022 & 1.0951 & fast jump \\
\bottomrule
\end{tabular}

\end{center}

\begin{center}
\begin{tikzpicture}
  \begin{groupplot}[
      group style={group size=3 by 1, horizontal sep=1.4cm}, % three panels
      width=5.6cm, height=6cm,                             % size of each panel
      grid=major, grid style={black!8},                    % faint grid
      axis background/.style={fill=black!2},               % very light background
      every axis plot/.append style={thick, mark size=1.2pt}, % same weight for all curves
    ]
    % ---- (i) solution error ----
    \nextgroupplot[
      xmode=log, ymode=log, title={(i) error at $t = 10$}, xlabel={$\Delta t$},
      legend to name=vdplegend, legend columns=7,          % legend below the panels
      legend style={font=\small, /tikz/every even column/.append style={column sep=0.25cm}},
    ]
    \addplot[m-order3, dashed, mark=triangle] table[x=dt, y=error] {data/error_rk3ws.dat};
    \addlegendentry{\texttt{rk3ws}}
    \addplot[m-rk4, mark=diamond*]            table[x=dt, y=error] {data/error_rk4.dat};
    \addlegendentry{\texttt{rk4}}
    \addplot[m-dp54, mark=pentagon*]          table[x=dt, y=error] {data/error_dp54.dat};
    \addlegendentry{\texttt{dp54}}
    % ---- (ii) period error ----
    \nextgroupplot[xmode=log, ymode=log, title={(ii) period error}, xlabel={$\Delta t$}]
    \addplot[m-order3, dashed, mark=triangle] table[x=dt, y=error] {data/period_rk3ws.dat};
    \addplot[m-rk4, mark=diamond*]            table[x=dt, y=error] {data/period_rk4.dat};
    \addplot[m-dp54, mark=pentagon*]          table[x=dt, y=error] {data/period_dp54.dat};
    % ---- (iii) stiffness ----
    \nextgroupplot[
      xmode=log, title={(iii) $\mu\,\Delta t_{\max}$}, xlabel={$\mu$},
      ymin=0.3, ymax=1.15, ytick={0.4,0.6,0.8,1.0},       % the scaled stable step
    ]
    \addplot[m-order3, dashed, mark=triangle] table[x=mu, y=scaled] {data/stiff_rk3ws.dat};
    \addplot[m-rk4, mark=diamond*]            table[x=mu, y=scaled] {data/stiff_rk4.dat};
    \addplot[m-dp54, mark=pentagon*]          table[x=mu, y=scaled] {data/stiff_dp54.dat};
    \addplot[m-order2, thin, densely dotted, domain=1:1000] {0.66667};   % h*/3, orders 1-2
    \addplot[m-order3, thin, densely dotted, domain=1:1000] {0.83758};   % h*/3, order 3
    \addplot[m-rk4, thin, densely dotted, domain=1:1000] {0.92843};      % h*/3, RK4
    \addplot[m-dp54, thin, densely dotted, domain=1:1000] {1.10219};     % h*/3, DP5(4)
  \end{groupplot}
  \node[below=0.2cm] at (current bounding box.south)        % legend centred under the panels
    {\pgfplotslegendfromname{vdplegend}};
\end{tikzpicture}
\end{center}
(i, ii) Every method converges at its general order in the solution
(\texttt{dp54} reaches round-off too early to measure, ---); in the period,
\texttt{rk3ws} shows order 3 instead of 2: it has effective order 3 (see the
nonlinear chemistry and Lorenz examples), its second-order error is a
near-identity change of variables $u \mapsto u + h^2X(u)$, and a change of
variables does not change the period. The period of the reference agrees to 13 digits with the
value in the literature, $6.6632868593\dots$ (iii) The scaled stable step
$\mu\,\Delta t_{\max}$ settles as $\mu$ grows, so the stable step falls like
$1/\mu$: the problem becomes stiff. Where it settles depends on where the
run first fails (last column, measured just above the limit). The methods
that fail on the slow branch (\texttt{rk3ws}) meet $h^*/3$ (dotted lines)
within $4\,\%$. The others fail
earlier, during a fast jump, where at $\mu\,\Delta t \approx 0.8$ a single
step spans most of the jump: a nonlinear failure that linear stability does
not predict (the estimate along the cycle is too optimistic for \texttt{rk4}
and \texttt{dp54}).
\end{document}
